\chapter{Model Validation and Hardware}\label{ch:validation}

\needswork{STATUS: DRAFTED.}

\section{Three independent implementations of the same physics}

The broadcast channel model is the load-bearing input to the capacity envelope, so it is validated
against two independent sources rather than one.

\begin{table}[h]
\centering
\caption{Saturation throughput (Mb/s), $1400$\,B broadcast frames. The slot-exact simulator was
written \emph{before} the published model was located and shares none of its assumptions.}
\label{tab:validation}
\begin{tabular}{rrrrrr}
\toprule
$N$ & published model & slot-exact sim & $\Delta$ & NS-3 & $\Delta$ \\
\midrule
5  & 4.3329 & 4.3305 & $-0.06\%$ & 4.3318 & $-0.03\%$ \\
10 & 3.1319 & 3.1321 & $+0.00\%$ & 3.1288 & $-0.10\%$ \\
20 & 1.7823 & 1.7830 & $+0.04\%$ & 1.7866 & $+0.24\%$ \\
35 & 1.2407 & 1.2407 & $-0.00\%$ & 1.2343 & $-0.51\%$ \\
50 & 1.2595 & 1.2585 & $-0.08\%$ & 1.2541 & $-0.43\%$ \\
\bottomrule
\end{tabular}
\end{table}

Agreement is $\leq 0.08\%$ against the independent simulator and $\leq 0.51\%$ against NS-3. Unicast
agrees with the acknowledged DCF model to $-0.40$/$+1.29\%$ across $N = 5$--$50$, with fixed-point
residuals below $10^{-12}$.

\begin{figure}[h]
\centering
\includegraphics[width=0.75\textwidth]{fig_ns3_bianchi.png}
\caption{Analytic models against simulation, both modes.}
\label{fig:validation}
\end{figure}

\section{The mechanism, confirmed by removing it}

The reason broadcast needs its own model is a specific asymmetry. A station that has just
transmitted redraws its backoff and may draw zero, taking the medium immediately after the
interframe space, while every station that deferred necessarily holds a counter of at least one. In
unicast the acknowledgement timeout blocks colliders from doing this; in broadcast nothing does.

The slot-exact simulator can disable \emph{exactly} that asymmetry and nothing else. Doing so
collapses it onto the naive reduction --- a $16.9\times$ gap at $N = 50$. \textbf{The mechanism
attribution is therefore not an inference from agreement; it is a controlled intervention.}

\begin{remark}[An expensive lesson]
This project used the naive reduction --- the unicast model with the acknowledgement removed and
the transmission probability fixed --- for several phases. It under-predicts by $16\times$ at
$N = 50$. The published broadcast model's own authors open by warning that unicast models cannot
simply be reduced for broadcast analysis. \emph{The failure was reading a model without reading its
preamble.}
\end{remark}

\section{A non-monotonicity that is real}

Saturation throughput is \textbf{not monotone decreasing} in $N$: it dips near $N \approx 35$ and
recovers. All three implementations reproduce it, so it is a property of the freeze-process series
rather than a bug --- each freeze stage contributes a hump peaking at a station count set by the
contention window.

This matters operationally. The capacity search breaks on the first violating $N$, which is valid
only if utilisation is monotone. It is --- because the explicit factor $N$ in the numerator outruns
the recovery --- but \textbf{that is an arithmetic accident, not a structural guarantee}, and it is
now verified exhaustively rather than assumed.

\section{Hardware}

Two Raspberry Pi 4 units in ad-hoc mode at $5$\,GHz measured broadcast link loss at
$p = 2.3 \times 10^{-4}$ and airtime at $1.995$\,ms per frame against $1.99$\,ms predicted --- an
independent check on the timing constants underlying every airtime figure in this thesis.

\begin{remark}[Two warnings, both earned]
This validates \emph{airtime}, not contention: one transmitter means zero contention, so the
broadcast contention model remains simulation-only.

And an initial $2.4$\,GHz sweep produced a tidy $97.45\%$ delivery figure that was \textbf{saturation
at the 802.11b $1$\,Mb/s broadcast basic rate}, not channel loss. It was caught by a pre-stated
prediction plus a load sweep, and is kept in the record labelled as what it was. Repeatability did
not protect against it --- the standard deviation across windows was small. \emph{A precisely
reproducible measurement of the wrong thing looks exactly like a good measurement.}
\end{remark}

\section{The utilisation ceiling, and a flaw in the author's own pre-registration}

The capacity envelope applies one measured utilisation ceiling to configurations whose frames span
$153$--$299$\,B, but that ceiling had only ever been measured at $288$\,B. A second sweep at
$174$\,B --- the compression-only baseline frame, which produces the denominators the capacity
ratios divide by --- tested it, with the prediction committed data-free beforehand.

\begin{center}
\begin{tabular}{lrr}
\toprule
& $288$\,B & $174$\,B \\
\midrule
$V \geq 0.95$ crossing & 2.435 & 2.367 \\
combined standard error & \multicolumn{2}{c}{$\pm 0.151$} \\
significance of the shift & \multicolumn{2}{c}{$0.45\sigma$} \\
\bottomrule
\end{tabular}
\end{center}

\textbf{Statistically indistinguishable across a $1.66\times$ change in frame size.} The universal
ceiling is justified and no published number moves.

\begin{remark}[The prediction that should not have been made]
The pre-registration also predicted a \emph{direction} --- that the crossing would drift higher ---
and offered a mechanism. It drifted lower, at $0.45\sigma$. The experiment has no power to resolve
the sign, so had the noise fallen the other way the author would have recorded a confirmation he
had not earned, and the mechanism offered is neither supported nor refuted.

This is recorded rather than quietly dropped because it is the same defect class the rest of this
thesis catalogues: \textbf{a number that looks like evidence and is not}. A directional prediction
needs a power estimate, or must be stated as a band. The load-bearing prediction did carry one,
which is why the run answers its question at all.
\end{remark}

\textbf{Still untested:} invariance in $N$. Both arms are at $N = 50$ while the envelope applies the
ceiling from $N = 2$ to $N = 213$. Frame-size invariance was the cheaper half of the assumption and
the half that varies across the envelope's rows; the other half is stated, not measured.
